% !PRINTED 584
\weberpage{584}
%<<CONT
equation (1), no linear relation with rational numerical coefficients can subsist between the $mn$ numbers (2).

Every number $\omega$ of the field $\Omega$ can be represented as an integral function of degree $(n-1)$ of $\Theta$ with coefficients in $R$, and every such number $\omega$ satisfies an equation irreducible in $R$ of degree at most $n$; $\omega$ is then and only then an integer, when this equation, after the coefficient of the highest power has been brought to 1, has integers in $R$ as coefficients.

The totality of all integers in $\Omega$ we denote, as before, by $\mathfrak{o}$.

If, without altering the numbers of the field $R$, we take for $\Theta$ all the roots $\Theta_1, \Theta_2, \ldots, \Theta_n$ of equation (1), we obtain $n$ fields
\weq{3}{\Omega_1 = R(\Theta_1), \ \Omega_2 = R(\Theta_2), \ \ldots, \ \Omega_n = R(\Theta_n),}
\gesperrt{which are conjugate with respect to the field $R$.}

If $\omega_1, \omega_2, \ldots, \omega_n$ are corresponding numbers of these fields, the product
\weq{4}{\mathfrak{N}_R(\omega) = \omega_1 \omega_2 \ldots \omega_n}
is called \gesperrt{the partial norm (relative norm) taken with respect to $R$} of the number $\omega$.

A number $\eta$ of the field $R$ has an ordinary norm $N_R(\eta)$, taken in this field, which is a rational number. Now the relative norm $\mathfrak{N}_R(\omega)$ is such a number $\eta$, and if we take its norm in the field $R$, we obtain the \gesperrt{total norm} of the number $\omega$ in the field of rational numbers:
\weq{5}{N_\Omega(\omega) = N_R\,\mathfrak{N}_R(\omega).}

This results if one expresses $\omega$ by the powers (2) of $\Theta$ and $\varrho$, and then puts for $\varrho$ the $m$ conjugate values, and for each of these $\varrho$ the values $\Theta$ conjugate with respect to $R$.

In the same sense the \gesperrt{functionals} of the field $\Omega$, and the \gesperrt{ideals} represented by these, now also have \gesperrt{their partial norms}. The partial norm of a functional in $\Omega$ is a functional in $R$, and accordingly the partial norm of an ideal in $\Omega$ is an ideal in $R$.

Of the conjugate ideals two or more can also be equal to one another, and indeed this can occur even when the conjugate functionals representing the ideals are not identical, provided only that they are associated. If one
%CONT>>
